% Seccion "Problem definition" del paper de bin packing 3D (BSG_BPP).
\section{Problem definition}
\label{sec:problem}

In the three-dimensional bin packing problem (3D-BPP) a set of $n$ rectangular boxes must be packed
into the minimum number of identical rectangular bins of length $L$, width $W$ and height $H$
\citep{Martello2000}. Box $i$ has dimensions $(l_i,w_i,h_i)$. A packing assigns every box to a bin
and to a position in it, such that every box lies entirely inside its bin, the boxes of a bin do not
overlap, and the faces of the boxes are parallel to those of the bin. In the typology of
\citet{Wascher2007} it is a three-dimensional single bin-size bin packing problem.

\paragraph{Orientation.}
In the classical version the orientation of the boxes is fixed: the dimensions $(l_i,w_i,h_i)$ are
parallel to $(L,W,H)$. In the version with rotations each box can be placed in any of its six
orientations. We consider both.

\paragraph{Vertical support.}
In practical loading every box above the floor must rest on the boxes below it. We use the criterion
that \citet{Elhedhli2019} take from industry practice: at least 70\% of the area of the base of every
box that is not on the floor must rest directly on the top faces of other boxes. Unlike
\citet{Elhedhli2019}, who form layers with a tolerance in the heights, we require exact contact, and we
impose the criterion on every box.
\pend{Check in \citet{Elhedhli2019} how the 70\% criterion and the 5~mm tolerance are applied, and
state the difference precisely.}

\paragraph{Box types.}
Many instances contain identical boxes. We group the boxes into $m$ types; type $t$ has dimensions
$(l_t,w_t,h_t)$, volume $v_t$ and $d_t$ boxes, so that $n=\sum_t d_t$. The methods of
Section~\ref{sec:method} only distinguish boxes by type.

\paragraph{Lower bounds.}
With fixed orientation we use the maximum of the continuous bound
$\lceil \sum_t d_t v_t/(LWH)\rceil$ and the bounds $L_1$ and $L_2$ of \citet{Martello2000}; with
rotations, the continuous bound. The lower bound, denoted $\mathrm{LB}$, is used only to stop the search
early and to report the gap of the solutions.
\pend{The implementation of $L_2$ admits in the volume term every box with $w\ge p$ and $h\ge q$, a
slightly larger set than in \citet{Martello2000}; add the one-line argument that the bound remains
valid (none of these boxes can be placed beside a box of the set $K_v$).}

\paragraph{Instances.}
The standard benchmark is the set of 320 instances generated by \citet{Martello2000}: eight classes,
$n\in\{50,100,150,200\}$ and ten instances per class and size. Most authors report totals over
classes 1 and 4--8 (240 instances) and over all eight classes. For vertical support we use the
industrial instances that accompany the benchmark of \citet{Elhedhli2019} (Section~\ref{sec:e-setup}).
